% !TEX root = ../qc-main.tex
\section*{学习地图}

\textbf{主要资料。}Michael A. Nielsen 与 Isaac L. Chuang，
\emph{Quantum Computation and Quantum Information}，
10th Anniversary Edition（十周年版），Cambridge University Press，2010。

\textbf{当前状态。}\S 1.1 笔记已确认；\S 1.2 四块讲解已完成，笔记于 2026-11-07 整理并获学习者确认。
两节均无独立编号 Example 或 Exercise，不设课程题目板块。下一学习节点为 \S 1.3 Quantum computation。

\begin{center}
\begin{tabularx}{\textwidth}{@{}lXll@{}}
\toprule
单元 & 内容 & 印刷页 / PDF 页 & 笔记状态 \\
\midrule
\hyperref[section:QC-C01-S01]{\S 1.1} & 总体视角 & 1--12 / 29--40 & 已确认 \\
\hyperref[section:QC-C01-S02]{\S 1.2} & 量子比特 & 13--17 / 41--45 & 已确认 \\
\bottomrule
\end{tabularx}
\end{center}

\textbf{当前单元。}\texttt{QC-C01-S02}：\S 1.2 Quantum bits，含 \S 1.2.1 Multiple qubits，
范围到印刷页 17 / PDF 页 45 的 \S 1.3 标题之前。四块内容如下。
\begin{center}
\begin{tabularx}{\textwidth}{@{}Xll@{}}
\toprule
知识块 & 印刷页 & PDF 页 \\
\midrule
\hyperref[topic:QC-C01-S02-T01]{单量子比特：状态、测量与物理实现} & 13--14 & 41--42 \\
\hyperref[topic:QC-C01-S02-T02]{相位、Bloch 球与读出信息} & 15--16 & 43--44 \\
\hyperref[topic:QC-C01-S02-T03]{多量子比特：张量积与局部测量} & 16 & 44 \\
\hyperref[topic:QC-C01-S02-T04]{Bell 态、纠缠关联与指数状态空间} & 16--17 & 44--45 \\
\bottomrule
\end{tabularx}
\end{center}

\textbf{先修衔接。}已学的线性代数、态矢量、幺正演化、张量积、自旋与 Pauli 矩阵足以支持入门。
密度算符与复合系统描述将系统复习，计算复杂度、编码与信息论按本书逐步补充；第 2、3 章分别提供量子力学和计算机科学基础。
